El objetivo principal de este Trabajo Fin de Grado es diseñar, implementar y validar una plataforma educativa de programación que integre un tutor basado en inteligencia artificial local con una metodología socrática de enseñanza y un entorno de ejecución de código aislado y seguro. Para alcanzar dicho objetivo, se plantean los siguientes objetivos específicos:

\begin{itemize}
    \item[--] Investigar el estado del arte en sistemas de tutoría inteligente y plataformas de aprendizaje de programación, analizando las capacidades actuales de la IA generativa para fomentar el pensamiento computacional.
    
    \item[--] Evaluar y seleccionar las soluciones tecnológicas más eficientes para la inferencia local de modelos de lenguaje (Ollama) y la orquestación de procesos aislados, garantizando un compromiso óptimo entre privacidad, rendimiento y seguridad.
    
    \item[--] Diseñar una arquitectura de sistema integral que armonice el procesamiento de lenguaje natural, un motor de ejecución de código efímero y una interfaz de usuario reactiva orientada a la experiencia pedagógica.
    
    \item[--] Adoptar una metodología de desarrollo iterativa basada en principios ágiles, permitiendo la evolución progresiva del sistema desde un Producto Mínimo Viable (MVP) hasta una plataforma robusta mediante ciclos de refinamiento continuo.
    
    \item[--] Elaborar la documentación técnica detallada siguiendo estándares de ingeniería del software, abarcando desde la captura de requisitos y casos de uso hasta el modelado de la arquitectura y diagramas estructurales.
    
    \item[--] Validar la plataforma mediante auditorías de seguridad sobre el entorno de ejecución y la realización de cuestionarios de experiencia de usuario para contrastar la efectividad real del tutor socrático.
\end{itemize}